\documentclass[10pt,a4paper]{amsart}
\usepackage[margin=1in]{geometry}
\usepackage{barnette-figs}
\usepackage{graphicx,amssymb}
\setlength{\parindent}{0pt}
\setlength{\parskip}{7pt}
\pagestyle{plain}
\begin{document}
\section*{04. The same argument with linear weights}
\begin{center}\begin{minipage}{.86\textwidth}\centering
\begin{tikzpicture}[bf,x=1.05cm,y=1.05cm]
\path[use as bounding box] (-5.80,-5.80) rectangle (5.80,5.80);
\fill[bfblack,even odd rule] (0,0) circle[radius=5.55] (0,0) circle[radius=5.0];
\fill[bfblack] (-4.330,-2.500) -- (-2.133,-0.405) -- (-2.200,3.528) -- (-4.330,-2.500) -- cycle;
\fill[bfblack] (-4.330,-2.500) -- (-0.344,-4.394) -- (-1.614,-2.366) -- (-4.330,-2.500) -- cycle;
\fill[bfblack] (-2.133,-0.405) -- (0.764,-0.167) -- (-0.963,1.243) -- (-2.133,-0.405) -- cycle;
\fill[bfblack] (-2.133,-0.405) -- (-1.614,-2.366) -- (0.900,-2.302) -- (-2.133,-0.405) -- cycle;
\fill[bfblack] (-0.344,-4.394) -- (4.330,-2.500) -- (0.900,-2.302) -- (-0.344,-4.394) -- cycle;
\fill[bfblack] (4.330,-2.500) -- (2.799,2.992) -- (2.506,-0.651) -- (4.330,-2.500) -- cycle;
\fill[bfblack] (0.000,5.000) -- (-2.200,3.528) -- (2.799,2.992) -- (0.000,5.000) -- cycle;
\fill[bfblack] (-2.200,3.528) -- (-0.963,1.243) -- (0.610,2.065) -- (-2.200,3.528) -- cycle;
\fill[bfblack] (0.764,-0.167) -- (2.799,2.992) -- (0.610,2.065) -- (0.764,-0.167) -- cycle;
\fill[bfblack] (0.764,-0.167) -- (0.900,-2.302) -- (2.506,-0.651) -- (0.764,-0.167) -- cycle;
\draw[bf edge] (0,0) circle[radius=5.0];
\draw[bf edge] (-4.330,-2.500) -- (-2.133,-0.405);
\draw[bf edge] (-4.330,-2.500) -- (-1.614,-2.366);
\draw[bf edge] (-4.330,-2.500) -- (-0.344,-4.394);
\draw[bf edge] (-4.330,-2.500) -- (-2.200,3.528);
\draw[bf edge] (-2.133,-0.405) -- (-1.614,-2.366);
\draw[bf edge] (-2.133,-0.405) -- (-2.200,3.528);
\draw[bf edge] (-2.133,-0.405) -- (-0.963,1.243);
\draw[bf edge] (-2.133,-0.405) -- (0.764,-0.167);
\draw[bf edge] (-2.133,-0.405) -- (0.900,-2.302);
\draw[bf edge] (-1.614,-2.366) -- (-0.344,-4.394);
\draw[bf edge] (-1.614,-2.366) -- (0.900,-2.302);
\draw[bf edge] (-0.344,-4.394) -- (4.330,-2.500);
\draw[bf edge] (-0.344,-4.394) -- (0.900,-2.302);
\draw[bf edge] (4.330,-2.500) -- (0.900,-2.302);
\draw[bf edge] (4.330,-2.500) -- (2.506,-0.651);
\draw[bf edge] (4.330,-2.500) -- (2.799,2.992);
\draw[bf edge] (0.000,5.000) -- (-2.200,3.528);
\draw[bf edge] (0.000,5.000) -- (2.799,2.992);
\draw[bf edge] (-2.200,3.528) -- (-0.963,1.243);
\draw[bf edge] (-2.200,3.528) -- (2.799,2.992);
\draw[bf edge] (-2.200,3.528) -- (0.610,2.065);
\draw[bf edge] (-0.963,1.243) -- (0.764,-0.167);
\draw[bf edge] (-0.963,1.243) -- (0.610,2.065);
\draw[bf edge] (0.764,-0.167) -- (0.900,-2.302);
\draw[bf edge] (0.764,-0.167) -- (2.506,-0.651);
\draw[bf edge] (0.764,-0.167) -- (2.799,2.992);
\draw[bf edge] (0.764,-0.167) -- (0.610,2.065);
\draw[bf edge] (0.900,-2.302) -- (2.506,-0.651);
\draw[bf edge] (2.506,-0.651) -- (2.799,2.992);
\draw[bf edge] (2.799,2.992) -- (0.610,2.065);
\begin{scope}\clip (-4.330,-2.500) -- (-2.133,-0.405) -- (-2.200,3.528) -- (-4.330,-2.500) -- cycle;\path[bf r] (-2.200,3.528) -- ++(-109.46:1.456) arc[start angle=-109.46,end angle=-89.02,radius=1.456] -- cycle;\end{scope}
\begin{scope}\clip (-4.330,-2.500) -- (-0.344,-4.394) -- (-1.614,-2.366) -- (-4.330,-2.500) -- cycle;\path[bf r] (-1.614,-2.366) -- ++(-177.18:0.603) arc[start angle=-177.18,end angle=-57.94,radius=0.603] -- cycle;\end{scope}
\begin{scope}\clip (-2.133,-0.405) -- (0.764,-0.167) -- (-0.963,1.243) -- (-2.133,-0.405) -- cycle;\path[bf r] (0.764,-0.167) -- ++(140.77:0.892) arc[start angle=140.77,end angle=184.70,radius=0.892] -- cycle;\end{scope}
\begin{scope}\clip (-2.133,-0.405) -- (-1.614,-2.366) -- (0.900,-2.302) -- (-2.133,-0.405) -- cycle;\path[bf r] (-2.133,-0.405) -- ++(-75.18:0.811) arc[start angle=-75.18,end angle=-32.02,radius=0.811] -- cycle;\end{scope}
\begin{scope}\clip (-0.344,-4.394) -- (4.330,-2.500) -- (0.900,-2.302) -- (-0.344,-4.394) -- cycle;\path[bf r] (-0.344,-4.394) -- ++(22.06:0.974) arc[start angle=22.06,end angle=59.26,radius=0.974] -- cycle;\end{scope}
\begin{scope}\clip (4.330,-2.500) -- (2.799,2.992) -- (2.506,-0.651) -- (4.330,-2.500) -- cycle;\path[bf r] (2.506,-0.651) -- ++(-45.39:0.576) arc[start angle=-45.39,end angle=85.40,radius=0.576] -- cycle;\end{scope}
\begin{scope}\clip (0.000,5.000) -- (-2.200,3.528) -- (2.799,2.992) -- (0.000,5.000) -- cycle;\path[bf r] (2.799,2.992) -- ++(144.34:1.211) arc[start angle=144.34,end angle=173.88,radius=1.211] -- cycle;\end{scope}
\begin{scope}\clip (-2.200,3.528) -- (-0.963,1.243) -- (0.610,2.065) -- (-2.200,3.528) -- cycle;\path[bf r] (-0.963,1.243) -- ++(27.59:0.691) arc[start angle=27.59,end angle=118.43,radius=0.691] -- cycle;\end{scope}
\begin{scope}\clip (0.764,-0.167) -- (2.799,2.992) -- (0.610,2.065) -- (0.764,-0.167) -- cycle;\path[bf r] (0.610,2.065) -- ++(-86.05:0.631) arc[start angle=-86.05,end angle=22.95,radius=0.631] -- cycle;\end{scope}
\begin{scope}\clip (0.764,-0.167) -- (0.900,-2.302) -- (2.506,-0.651) -- (0.764,-0.167) -- cycle;\path[bf r] (0.900,-2.302) -- ++(45.79:0.856) arc[start angle=45.79,end angle=93.64,radius=0.856] -- cycle;\end{scope}
\begin{scope}\clip (-4.330,-2.500) -- (-2.133,-0.405) -- (-2.200,3.528) -- (-4.330,-2.500) -- cycle;\path[bf s] (-2.133,-0.405) -- ++(90.98:0.572) arc[start angle=90.98,end angle=223.64,radius=0.572] -- cycle;\end{scope}
\begin{scope}\clip (-4.330,-2.500) -- (-0.344,-4.394) -- (-1.614,-2.366) -- (-4.330,-2.500) -- cycle;\path[bf s] (-0.344,-4.394) -- ++(122.06:0.957) arc[start angle=122.06,end angle=154.58,radius=0.957] -- cycle;\end{scope}
\begin{scope}\clip (-2.133,-0.405) -- (0.764,-0.167) -- (-0.963,1.243) -- (-2.133,-0.405) -- cycle;\path[bf s] (-0.963,1.243) -- ++(-125.37:0.709) arc[start angle=-125.37,end angle=-39.23,radius=0.709] -- cycle;\end{scope}
\begin{scope}\clip (-2.133,-0.405) -- (-1.614,-2.366) -- (0.900,-2.302) -- (-2.133,-0.405) -- cycle;\path[bf s] (-1.614,-2.366) -- ++(1.46:0.648) arc[start angle=1.46,end angle=104.82,radius=0.648] -- cycle;\end{scope}
\begin{scope}\clip (-0.344,-4.394) -- (4.330,-2.500) -- (0.900,-2.302) -- (-0.344,-4.394) -- cycle;\path[bf s] (0.900,-2.302) -- ++(-120.74:0.608) arc[start angle=-120.74,end angle=-3.30,radius=0.608] -- cycle;\end{scope}
\begin{scope}\clip (4.330,-2.500) -- (2.799,2.992) -- (2.506,-0.651) -- (4.330,-2.500) -- cycle;\path[bf s] (2.799,2.992) -- ++(-94.60:1.462) arc[start angle=-94.60,end angle=-74.42,radius=1.462] -- cycle;\end{scope}
\begin{scope}\clip (0.000,5.000) -- (-2.200,3.528) -- (2.799,2.992) -- (0.000,5.000) -- cycle;\path[bf s] (-2.200,3.528) -- ++(-6.12:1.042) arc[start angle=-6.12,end angle=33.79,radius=1.042] -- cycle;\end{scope}
\begin{scope}\clip (-2.200,3.528) -- (-0.963,1.243) -- (0.610,2.065) -- (-2.200,3.528) -- cycle;\path[bf s] (0.610,2.065) -- ++(152.50:0.710) arc[start angle=152.50,end angle=207.59,radius=0.710] -- cycle;\end{scope}
\begin{scope}\clip (0.764,-0.167) -- (2.799,2.992) -- (0.610,2.065) -- (0.764,-0.167) -- cycle;\path[bf s] (0.764,-0.167) -- ++(57.21:0.895) arc[start angle=57.21,end angle=93.95,radius=0.895] -- cycle;\end{scope}
\begin{scope}\clip (0.764,-0.167) -- (0.900,-2.302) -- (2.506,-0.651) -- (0.764,-0.167) -- cycle;\path[bf s] (2.506,-0.651) -- ++(164.47:0.723) arc[start angle=164.47,end angle=225.79,radius=0.723] -- cycle;\end{scope}
\draw[bf D] (-2.200,3.528) -- (-2.133,-0.405);
\draw[bf D] (-1.614,-2.366) -- (-0.344,-4.394);
\draw[bf D] (0.764,-0.167) -- (-0.963,1.243);
\draw[bf D] (-2.133,-0.405) -- (-1.614,-2.366);
\draw[bf D] (-0.344,-4.394) -- (0.900,-2.302);
\draw[bf D] (2.506,-0.651) -- (2.799,2.992);
\draw[bf D] (2.799,2.992) -- (-2.200,3.528);
\draw[bf D] (-0.963,1.243) -- (0.610,2.065);
\draw[bf D] (0.610,2.065) -- (0.764,-0.167);
\draw[bf D] (0.900,-2.302) -- (2.506,-0.651);
\draw[bf root] (-4.330,-2.500) circle[radius=2.3pt];
\fill[bf vertex] (-2.133,-0.405) circle[radius=1.4pt];
\fill[bf vertex] (-1.614,-2.366) circle[radius=1.4pt];
\fill[bf vertex] (-0.344,-4.394) circle[radius=1.4pt];
\draw[bf root] (4.330,-2.500) circle[radius=2.3pt];
\draw[bf root] (0.000,5.000) circle[radius=2.3pt];
\fill[bf vertex] (-2.200,3.528) circle[radius=1.4pt];
\fill[bf vertex] (-0.963,1.243) circle[radius=1.4pt];
\fill[bf vertex] (0.764,-0.167) circle[radius=1.4pt];
\fill[bf vertex] (0.900,-2.302) circle[radius=1.4pt];
\fill[bf vertex] (2.506,-0.651) circle[radius=1.4pt];
\fill[bf vertex] (2.799,2.992) circle[radius=1.4pt];
\fill[bf vertex] (0.610,2.065) circle[radius=1.4pt];
\end{tikzpicture}\par $F(1+xa)=9x^2-\frac{2563}{4056}x^3+O(x^4)$, of degree ten.\par Exact coefficients are recorded in computed.json.
\end{minipage}\end{center}
Replace each exponential spoke weight by one plus x times a rational circulation weight. Multiply the selected weights for each second state and multiply by the product of local edge signs. Summing over the eight pairs gives the polynomial below the picture. Its coefficients are rational after this substitution; the polynomial in independent spoke variables has integer coefficients. Fixed-state cancellation still removes cyclic states because each cancelled pair selects exactly the same monomial. In the other grouping, each D-cycle contributes a difference of products whose first coefficient is positive white area. All groups with minimum cycle count have the same sign. Therefore the first surviving coefficient cannot cancel. The displayed initial expansion checks this mechanism; the full coefficients are saved alongside it. It does not remove the geometric input: positive circulation is still the reason nonvanishing holds.

\textbf{Status.} PROVED-HERE in the determinant referee, SIGNIFICANCE Section 2; independently AI-refereed reformulation, human review pending. COMPUTED polynomial here.

\textbf{Connects to.} 2 exponential terms; 7 white surplus; 10 determinant; 25 Berezin.

{\footnotesize\raggedright Sources: research/\allowbreak{}2026-11-14\_fable\_referee-determinants/\allowbreak{}SIGNIFICANCE.md, Section 2;\allowbreak{} build.py;\allowbreak{} computed.json.\par}

\end{document}